\scanpage{136}
\mspage{-III- \handwritten{2} -}
\begin{flushright}n°\handwritten{383}\end{flushright}
Considérons les \uline{chambres} définies par $\underline{H}$, c'est-à-dire les composantes connexes de la réunion des hyperplans $H\in\underline{H}$ (\S\ 1, n°1) ; rappelons qu'une chambre est un cône ouvert convexe (ibid.).

PROPOSITION 1.- \uline{Soient} $C$ \uline{une chambre}, $H_1,\ldots,H_m$ \uline{ses murs}. \uline{Les réflexions} $s_{H_1},\ldots,s_{H_m}$ \uline{engendrent} $G$. \uline{Le groupe} $G$ \uline{opère transitivement sur l'ensemble des chambres définies par} $\underline{H}$.

Soient $G'$ le sous-groupe de $G$ engendré par $s_{H_1},\ldots,s_{H_m}$, $x$ un point de $C$, $C'$ une autre chambre, et $x'$ un point de $C'$. Parmi tous les transformés $g'(x')$ $(g'\in G')$, soit $y$ l'un de ceux qui sont les plus voisins de $x$. Si $y$ n'était pas dans $C$, il existerait un mur $H_i$ de $C$ tel que $x$ et $y$ soient de part et d'autre de $H_i$ ; alors le symétrique $s_{H_i}(y)$ serait strictement plus proche de $x$ que $y$ ; on a donc $y\in C$, ce qui montre que $G'$ opère transitivement sur l'ensemble des chambres définies par $\underline{H}$, et prouve la seconde assertion.

Soit maintenant $H\in\underline{H}$. D'après le \S\ 1, n°1, lemme 3, $H$ est un mur d'une chambre $C'$ définie par $\underline{H}$. Il existe donc, d'après ce qu'on vient de voir, un élément $g'$ de $G'$ et un indice $i$ $(1\handwritten{\leq}i\handwritten{\leq}m)$ tels que $H=g'(H_i)$. D'après (1) on a donc $s_H=g's_{H_i}{g'}^{-1}$, d'où $s_H\in G'$ ; ceci montre que $G=G'$. C.Q.F.D.

Nous fixerons désormais une chambre $C$ (définie par $\underline{H}$), noterons $H_1,\ldots,H_m$ ses murs, et poserons $s_{H_i}=s_i$. Soit $e_i$ le vecteur unitaire orthogonal à $H_i$ et situé du même côté de $H_i$ que $C$. La chambre $C$ est donc l'ensemble des points $x$ de $E$ tels que
\[
\begin{array}{l@{\qquad}l}(2)&(x\handwritten{|}e_i)\handwritten{>}0\qquad\text{pour tout }i=1,\ldots,m.\end{array}
\]

PROPOSITION 2.- \uline{Soient} $i,j$ \uline{deux indices distincts}, \uline{et} $m_{ij}$ \uline{l'ordre de} $s_i s_j$ \uline{dans} $G$. \uline{On a alors} $(e_i\handwritten{|}e_j)=-\cos(\handwritten{\pi}/m_{ij})$, \uline{de sorte que} $(e_i\handwritten{|}e_j)\handwritten{\leq}0$. \uline{Si} $(e_i\handwritten{|}e_j)\handwritten{<}0$, \uline{on a} $(e_i\handwritten{|}e_j)\handwritten{\leq}-1/2$.

Soit $G^{\handwritten{\unsure{o|0|\circ}}}$ le sous-groupe de $G$ engendré par $s_i$ et $s_j$ ; le sous-espace essentiel de $E$
% archive sheet number: 135
